\section{Hybride Darstellung und Blockstruktur}
\subsection{Motivation}
Rein parametrische Modelle sind niedrigdimensional und interpretierbar, aber
für reale Daten wenig flexibel. Rein abgetastete Modelle sind maximal flexibel,
führen jedoch zu hochdimensionalen, oft schlecht konditionierten Problemen.
Eine hybride Formulierung verbindet beide Vorteile.

\subsection{Hybride Darstellung}
\begin{definition}[Hybrides Krümmungsmodell]
\[
\kappa(s)=\kappa_{\mathrm{param}}(s;\mathbf T)+\delta\kappa(s),
\]
wobei \(\kappa_{\mathrm{param}}\) ein strukturiertes parametrisches Modell und
\(\delta\kappa(s)\) ein Korrekturterm ist.
\end{definition}
\begin{definition}[Hybrides Überhöhungsmodell]
\[
\phi(s)=\phi_{\mathrm{param}}(s;\mathbf T_\phi)+\delta\phi(s).
\]
\end{definition}

\subsection{Diskretisierung der Korrekturen}
\[
\delta\kappa_i=\delta\kappa(s_i),\qquad
\delta\phi_i=\delta\phi(s_i).
\]
\begin{definition}[Erweiterter Parametervektor]
\[
\mathbf x=(\mathbf T,\mathbf T_\phi,
\delta\boldsymbol\kappa,\delta\boldsymbol\phi).
\]
\end{definition}

\subsection{Regularisierung}
Zur Vermeidung von Überanpassung gilt
\[
J_\delta=\int_0^L(\delta\kappa'(s))^2\,\dd s
+\int_0^L(\delta\phi'(s))^2\,\dd s.
\]
Damit erfasst die Korrektur nur glatte, physisch bedeutsame Abweichungen.

\subsection{Residuenstruktur}
Der Residuenvektor besteht aus parametrischen Beiträgen,
Korrekturbeiträgen, physischen Bedingungen und Datenanpassungstermen.

\subsection{Blockstruktur der Variablen}
Der Parametervektor zerfällt natürlich in
\[
(\mathbf T,\delta\boldsymbol\kappa)
\quad\text{und}\quad
(\mathbf T_\phi,\delta\boldsymbol\phi).
\]
Es entsteht eine zweistufige Struktur mit parametrischen Variablen im groben
und Korrekturvariablen im feinen Maßstab.

\subsection{Blockstruktur der Jacobi-Matrix}
\begin{definition}[Block-Jacobi-Matrix]
\[
J_r=\begin{pmatrix}
J_{\kappa\kappa}^{\mathrm{param}}&J_{\kappa\kappa}^{\mathrm{corr}}&0&0\\
0&0&J_{\phi\phi}^{\mathrm{param}}&J_{\phi\phi}^{\mathrm{corr}}\\
J_{\mathrm{dyn},\kappa}&J_{\mathrm{dyn},\kappa}^{\delta}&
J_{\mathrm{dyn},\phi}&J_{\mathrm{dyn},\phi}^{\delta}
\end{pmatrix}.
\]
\end{definition}
Parametrische und Korrekturterme sind schwach gekoppelt; Krümmung und
Überhöhung weitgehend separierbar; die Kopplung erfolgt vor allem durch die
Dynamik.

\subsection{Hierarchische Interpretation}
Die parametrische Ebene definiert die globale Alignment-Struktur, die
Korrekturebene verfeinert lokale Abweichungen. Dies entspricht der Praxis:
Entwurf definiert Absicht, Bau führt Abweichungen ein.

\subsection{Rechenstrategie}
Zunächst werden parametrische Variablen optimiert, dann Korrekturterme
verfeinert; optional wechseln beide Ebenen.

\subsection{Verbindung zum Detaillierungsgrad}
Grober LOD verwendet nur das parametrische Modell, feiner LOD Parametrik und
Korrekturen.

\subsection{Zentrale Erkenntnis}
\begin{proposition}
Die hybride Formulierung trennt strukturelle Modellierung von Datenanpassung
und bewahrt zugleich ein dünn besetztes, gut konditioniertes Problem.
\end{proposition}

\subsection{Verbindung zu AIM}
\begin{summary}
Hybride Modellierung und Blockstruktur ermöglichen interpretierbare
Parametrik, robuste Behandlung unvollkommener Daten und effiziente Optimierung
durch dünne Besetzung und Lokalität. Sie verbinden Engineering-Entwurf mit
datengetriebener Rekonstruktion.
\end{summary}
